\chapter{統計力学と量子統計 --- 個性の喪失とマクロな奇跡、そしてAIとの交差点}

\section*{本章の学習目標}
\begin{itemize}
\item 大数の法則とエルゴード仮説から、ミクロの「確率」がいかにしてマクロの「確実性」に変わるかを知る。
\item ミクロとマクロを繋ぐ「統計力学」の心臓部である「分配関数」の概念を理解する。
\item 古典力学と量子力学の決定的な違いである「同種粒子の識別不可能性（個性の喪失）」を知る。
\item フェルミ・ディラック統計：パウリの排他原理がもたらす「電子の満員電車」と、星の崩壊を防ぐ力を学ぶ。
\item ボース・アインシュタイン統計：粒子が一斉に同期する「ボース・アインシュタイン凝縮（BEC）」の奇跡を知る。
\item 統計力学のイジングモデルやボルツマン分布が、AIの「ボルツマンマシン」や「ソフトマックス関数」と全く同じ数学的構造を持っていることを理解する。
\end{itemize}

\section{ミクロからマクロへ：「大数の法則」と統計力学の完成}

第10章において、私たちは「シュレディンガー方程式」という、ミクロな粒子1個の振る舞いを記述する究極のルールを手に入れました。 しかし、私たちが日常で手にするコップ1杯の水の中には、約 \ensuremath{10^{23}} 個（1000垓個：アボガドロ定数）という途方もない数の水分子が含まれています。これらすべての分子について方程式を立ててスーパーコンピュータで解こうとすれば、「次元の呪い」により、宇宙の寿命が尽きても計算はほとんど進みません。

そこで、第3章（熱力学）で登場したルートヴィッヒ・ボルツマンや、アメリカの物理学者ウィラード・ギブズらが完成させた\textbf{「統計力学（Statistical Mechanics）」}の出番となります。 統計力学は、粒子の1個1個の動きを追跡するのを潔く諦め、「サイコロの確率（統計）」を使って、\ensuremath{10^{23}} 個の集団全体が持つマクロな性質（温度、圧力、エントロピー）を計算する理論です。

\subsection{1000垓個のサイコロと「揺らぎの消失」}

統計力学の根底にあるのは、確率論における\textbf{「大数の法則（Law of Large Numbers）」}という数学的マジックです。

例えば、コインを10回投げたとき、表が8回（80\%）出ることは珍しくありません。期待値の50\%からの「ズレ（揺らぎ）」が大きい状態です。しかし、コインを100万回、1兆回と投げれば、表が出る確率は限りなく「50\%」にピタリと落ち着きます。 数学（統計学）の定理によれば、この「相対的な揺らぎ（平均値に対するズレの割合）」の大きさは、粒子数 \(N\) の平方根に反比例します。

\[
\text{相対的な揺らぎ} \sim \frac{1}{\sqrt{N}}
\]

これをコップの水分子に当てはめてみましょう。粒子数 \(N\) は \ensuremath{10^{23}} ですから、その揺らぎの大きさは \ensuremath{1/\sqrt{10^{23}}} \ensuremath{\approx} \ensuremath{10^{-11.5}}（約1000億分の1）となります。 個々の分子はデタラメ（カオス）に飛び回り、シュレディンガー方程式に従って激しく揺らいでいます。しかし、数が \ensuremath{10^{23}} 個にもなると、その個別の「揺らぎ」は互いにほとんど打ち消し合い、マクロにはほぼブレない安定した量へと姿を変えるのです。

さらに言えば、「温度」や「圧力」といった概念は、水分子「1個」には存在しません。ミクロな偶然が途方もない数集まることで、初めてマクロな必然（安定した物理量）へと\textbf{「創発」}するのです。これが統計力学の第一の見どころです。

\subsection{システムとアンサンブル：現実と「並行世界」の対比}

統計力学を理解する上で、非常に重要となる2つの概念があります。それが\textbf{「システム（系）」と「アンサンブル（統計集団）」}です。

システム（系：System）＝ 現実の1つの対象 物理学において、私たちが観測や計算の対象として宇宙全体から境界線を引いて切り出した「現実の物質のまとまり」のことです。例えば「目の前にあるコップ1杯の水」や「シリンダーに閉じ込めた気体」などです。 現実の温度や圧力を計算するには、この1つのシステムの中の分子が時間をかけて色々な状態を巡るのを、ストップウォッチを持ってじっと見守り、平均をとる（時間平均）しかありません。

アンサンブル（統計集団：Ensemble）＝ 架空の無数のコピー しかし、分子がどう動くかを長時間待って計算するのは不可能です。そこでギブズは、確率計算を数学的に可能にするための極めてエレガントな思考実験を行いました。それが「アンサンブル」です。 これは、目の前のシステムと\textbf{「全く同じマクロな条件（同じ温度、同じ体積など）を持った、無限個のコピーがズラリと並ぶ『並行世界』の集まり」}を頭の中で想像したものです。そして、ある一瞬にこの無限のコップを「パシャッ」と一斉に写真に撮り、そのすべての状態を足し合わせて平均をとる（アンサンブル平均）という数学モデルに置き換えました。

ここでボルツマンが提唱した\textbf{「エルゴード仮説（Ergodic Hypothesis）」が登場します。これは「水分子たちは、十分な時間が経てば、エネルギーが許す限りの『あらゆる状態の組み合わせ（配置パターン）』をすべて平等に巡り尽くすはずだ」という大前提です。 この仮説を認めれば、「1つの現実のシステムを無限の時間をかけて測った平均（時間平均）」と、「無限個の架空のコピーを一瞬で測った平均（アンサンブル平均）」は数学的に完全に一致}することになります。 （※実はこのエルゴード仮説は、現代の数学や物理学をもってしても「すべての物質で本当に成り立つのか」が完全には証明されていない、非常に奥深いミステリーの一つです。しかし、現実の気体や液体では驚くほど見事に成り立つため、統計力学の強固な土台となっています。）

このアンサンブルの導入は、物理学の計算から「長時間待つ」という厄介な手続きを取り除き、純粋な「あり得る状態の確率の足し算（積分）」によって現実のシステムを予測する道を開きました。 （※ちなみにこの「アンサンブル」という言葉と思想は、現代のAIにおいても、複数の異なるモデルの予測を掛け合わせてブレをなくす「アンサンブル学習（ランダムフォレストなど）」や、気象予測のカオスを扱う「アンサンブル予報」として、同じ発想で引き継がれています）。

\subsection{統計力学のマスターキー：「分配関数（Partition Function）」}

では、その「並行世界（アンサンブル）」において、ある特定の状態が観察される確率はどのように計算されるのでしょうか？ 統計力学において、システム全体がある特定のエネルギー状態（\(E_i\)）をとる確率は、温度（\(T\)）によって決まり、次のような\textbf{「ボルツマン分布」}に従うことが証明されています。

\[
P_i = \frac{e^{-E_i/k_B T}}{Z}
\]

（\(k_B\) はボルツマン定数。エネルギー \(E\) が低い状態ほど確率が高く、温度 \(T\) が上がると高いエネルギー状態の確率も増えることを意味します）

そして、この確率を正確に計算するために「あり得るすべての状態の確率の足し算（分母）」をまとめたものを、\textbf{「分配関数（Partition Function：\(Z\)）」}と呼びます。

\[
Z = \sum_i e^{-E_i/k_B T}
\]

分配関数 \(Z\) は、一見するとただの分母（全確率を1にするための規格化定数）に過ぎません。しかし物理学者にとって、この \(Z\) はあらゆる願いを叶えてくれる\textbf{「統計力学のマスターキー（魔法の箱）」}なのです。

なぜなら、この \(Z\) さえ計算できれば、その対数（\(\ln Z\)）をとるだけで、第3章で学んだ「システムが外に仕事を取り出せる有効なエネルギー」である\textbf{「ヘルムホルツの自由エネルギー（\(F\)）」}が次のように一発で導き出されるからです。

\[
F = -k_B T \ln Z
\]

そして、この自由エネルギー \(F\) を手に入れれば、あとはそれを「温度（\(T\)）」で微分すれば『エントロピー（\(S\)）』が出てきます。「体積（\(V\)）」で微分すれば『圧力（\(P\)）』が出てきます。 つまり、ミクロな原子のエネルギー状態をすべて足し合わせた「分配関数 \(Z\)」を組み立てることさえできれば、あとは微分を繰り返すだけで、その物質が持つマクロな性質（温度、圧力、比熱、エントロピーなど）を芋づる式に導き出せるのです。

このようにして完成した統計力学は、マクロ（熱力学）とミクロ（力学）を完璧な数学で繋ぐ架け橋となりました。しかし、次節で見るように、20世紀に入って「量子力学」が誕生したことで、この確率の計算ルールは「粒子の個性の喪失」という強烈な洗礼を受け、さらなる奇跡の現象を生み出すことになります。

\section{量子力学の洗礼：粒子の「個性」が消滅する世界}

古典力学では、同じ種類の粒子であっても「粒子A」「粒子B」とラベルを貼って区別できると考えます。たとえば2個のボールが左右の箱に1個ずつ入っているとき、「Aが左でBが右」と「Bが左でAが右」は、ボールに名前を書いておけば別々の状態として数えられます。

しかし、量子力学のミクロな世界（電子や光子など）では、状況が全く異なります。 宇宙に存在する「電子」は、どれを持ってきても質量や電荷が1ミリの狂いもなく完全に同一です（傷一つありません）。 さらに決定的なのが、第10章で学んだ\textbf{「不確定性原理」と「波としての性質」}です。量子力学では、粒子の位置とスピードを同時に正確に測ることはできず、電子は「確率の波」として空間に広がっています。

もし2つの電子が近づいてきて衝突し、弾け飛んだとします。このとき、波と波が重なり合って（干渉して）しまうため、「どちらの電子がどっちへ飛んでいったか」という軌道を追跡することは原理的に不可能です。 つまり、量子力学において同種の粒子（電子同士、光子同士など）は、\textbf{「原理的に識別不可能（個性が完全に存在しない）」}のです。電子Aと電子Bを入れ替えても、宇宙から見れば「全く同じ1つの状態」にしかなりません。

\subsection{確率の激変と「分配関数」の書き換え}

この「個性の喪失（識別不可能性）」は、先ほどの確率の計算（場合の数）を根底から狂わせます。 【量子力学の計算例】 区別がつかない2つの電子を、「右の箱」と「左の箱」にそれぞれ1個ずつ入れるとします。

パターン：左に1個、右に1個

これしかありません。AとBの入れ替えが意味を持たないため、答えは\textbf{「1通り」}になってしまうのです。

古典力学では「異なる2通りの状態」として数えていたものが、量子力学では「1通りの状態」としてカウントされてしまう。 第3章で学んだように、エントロピーやマクロな物理法則は「パターンの数（場合の数 \(W\)）」によって決まります（\(S = k_B \log W\)）。パターンの数が半分に減ってしまうということは、物質の熱の伝わり方や圧力の計算が、古典力学の予測から劇的にズレてしまうことを意味します。

この「粒子に個性が存在しないこと」を前提にして確率（分配関数）の計算ルールを根底から書き直した理論を\textbf{「量子統計」と呼びます。 そして驚くべきことに、宇宙の粒子はその生まれ持ったスピン（自転のような性質）の違いによって、この量子統計のルールの下でもさらに「2つの全く異なるグループ」}に分かれることが判明したのです。

\section{フェルミ・ディラック統計：宇宙の崩壊を防ぐ「満員電車」}

グループの一つ目は、電子、陽子、中性子、クォークなど、物質の「材料」となる粒子たちです。これらを\textbf{「フェルミ粒子（Fermion）」}と呼びます。イタリアのエンリコ・フェルミと、イギリスのポール・ディラックがその統計法則を導き出しました。

\subsection{パウリの排他原理と波動関数の「反対称性」}

フェルミ粒子には、量子力学における絶対的なルール\textbf{「パウリの排他原理」が適用されます。1925年にオーストリアの物理学者ヴォルフガング・パウリが提唱したこの原理は、「一つの状態（エネルギーの席やスピンの向きが完全に同じ状態）には、絶対に1つの粒子しか座ることができない」}という、極めて強い反発のルールです。

なぜこのようなルールがあるのでしょうか？ 第10章で学んだ波動関数の言葉で説明すると、フェルミ粒子の波には\textbf{「2つの粒子の位置を入れ替えると、波のプラスとマイナス（符号）が完全にひっくり返る（反対称性）」}という奇妙な数学的性質があります。

もし2つの電子が「全く同じ場所、全く同じ状態」にあったとします。その2つを入れ替えても、同じ状態なのだから波の形は変わらないはずです。しかし、フェルミ粒子のルールでは「入れ替えたら符号がマイナスにならなければならない」のです。 「自分自身と、自分のマイナスが等しい（X = -X）」という数字は、数学的には「ゼロ」しかありません。つまり、\textbf{「全く同じ状態の電子が2つ重なると、波が打ち消し合って存在確率がゼロになってしまう（存在できない）」}のです。これが排他原理の数学的な正体です。

\subsection{フェルミ・ディラック分布関数：確率の「上限」を決める数式}

このパウリの排他原理を組み込んで、あるエネルギー（\(E\)）の席に粒子が座る確率（平均粒子数）を計算したものが、次の\textbf{「フェルミ・ディラック分布関数」}です。

\[
f(E) = \frac{1}{e^{(E-\mu)/k_B T}+1}
\]

（\(f(E)\) は粒子が座る確率、\(\mu\) は化学ポテンシャル、\(k_B\) はボルツマン定数、\(T\) は絶対温度）

ここで登場する \(\mu\)（化学ポテンシャル）とは、本来は熱力学で「システムに粒子を1つ追加するために必要なエネルギー」を意味する言葉ですが、この数式においては\textbf{「席がちょうど50\%の確率で埋まる境界線（水面の高さ）」}を表す基準値として働きます。

数式を見てください。もしエネルギー \(E\) が \(\mu\) と全く同じ値（\(E = \mu\)）だとすると、分母の指数部分は \(e^{0} = 1\) となり、全体の確率は \(\frac{1}{1+1} = \frac{1}{2}\)（50\%）になります。 絶対零度（熱の揺らぎが全くない状態）のとき、この \(\mu\)（フェルミエネルギーと呼ばれます）を境にして、それよりエネルギーが低い席は確率100\%で完全に満席となり、高い席は確率0\%で完全に空席になります。まさに「水面」のように、電子が下からギッシリと積み重なっている様子を、この数式は見事に表現しているのです。

11.1節で見た古典的なボルツマン分布の分母に、たった一つ\textbf{「\(+1\)」という数字が付け足されています。 しかし、この「\(+1\)」こそが、フェルミ粒子の満員電車ルールを数式上で表しています。分母に必ず \(1\) が足されるため、全体の確率 \(f(E)\) は絶対に「\(1\)（100\%）」を超えることができません}。どんなに温度が低くても、どんなにエネルギー状態が低くても、「1つの席には最大1個までしか座れない」という事実が、このシンプルな \(+1\) によって数学的に保証されているのです。

\subsection{フェルミの海：ペチャンコにならない私たちの体}

温度が絶対零度（マイナス273度）まで冷えたときを想像してください。 古典力学の粒子なら、全員がエネルギーゼロの「一番下の席」にドッと落ち込んで重なり合います。しかしフェルミ粒子は、一つの席に一人しか座れないため、一番下の席が埋まると、次の粒子は一つ上の席、その次はさらに上の席\ldots{}\ldots{}と、下から順番にギッシリと積み重なっていきます。これを\textbf{「フェルミの海（満員電車状態）」}と呼びます。

この性質こそが、私たちの体がペチャンコに潰れない理由です。電子がフェルミ粒子であるおかげで、原子の周りの電子は一番下の軌道にすべて落ち込むことができず、外側へ外側へと空間を占有します。物質に「体積（大きさ）」や「硬さ」があるのは、このフェルミ粒子の「席取りゲームの積み重ね」のおかげなのです。

さらに、寿命を迎えた星が重力で潰れようとするのをギリギリで支えているのも、このフェルミ粒子の反発力（縮退圧）です。「白色矮星」では電子の縮退圧が、「中性子星」では中性子の縮退圧が重力に抵抗します。ただし、星の質量が十分に大きい場合には縮退圧でも支えきれず、最終的にブラックホールへ崩壊します。排他原理は宇宙の崩壊を何でも止める万能の壁ではなく、一定の範囲で重力に対抗する量子力学的な圧力なのです。

\section{ボース・アインシュタイン統計：一斉に同期する「マクロな波」}

グループの二つ目は、光子（フォトン）や、力を伝える粒子たちです。これらを\textbf{「ボース粒子（Boson）」}と呼びます。インドの物理学者サティエンドラ・ボースの論文をアインシュタインが拡張したことで導き出されました。

\subsection{波動関数の「対称性」と群れたがる性質}

ボース粒子はフェルミ粒子とは真逆で、パウリの排他原理に一切縛られません。 11.3節で「フェルミ粒子の波は入れ替えるとマイナスになる（反対称性）」と説明しましたが、ボース粒子の波は\textbf{「2つの粒子の位置を入れ替えても、波の形が全く変わらない（対称性）」}という数学的性質を持っています。入れ替えてもマイナスにならないため、同じ状態に何個重なっても波が打ち消し合うことがありません。

それどころか、ボース粒子には\textbf{「一つの状態（席）に、無限に何個でも重なって入りたがる（みんなと同じ状態になりたがる）」}という極めて社交的な性質があります。 古典的な確率論なら、空いている席と人が座っている席があれば、人はランダムに散らばるはずです。しかしボース粒子の確率計算（ボース・アインシュタイン統計）では、「すでに粒子がたくさん入っている席ほど、次の粒子がさらにそこに入りやすくなる」という奇妙なエコシステムが働きます。まるで流行の店に行列がさらに人を呼ぶように、ボース粒子は同じ状態へと猛烈な勢いで群れていくのです。

\subsection{ボース・アインシュタイン分布関数：無限の「群れ」を許す数式}

この群れたがる性質を数式で表したのが、次の\textbf{「ボース・アインシュタイン分布関数」}です。

\[
f(E) = \frac{1}{e^{(E-\mu)/k_B T}-1}
\]

フェルミ粒子の式にあった分母の「\(+1\)」が、ここでは\textbf{「\(-1\)」に変わっています。 たったこれだけの違いですが、結果は天地ほど異なります。分母が引き算（\(-1\)）になるため、ある条件（\(E\) が \(\mu\) に近づく状態）では、分母が限りなく「\(0\)」に近づき、結果として \(f(E)\) の値は「無限大（\(\infty\)）」に向かって爆発}することができます。つまり、「1つの席に何億個、何兆個という粒子が無限に重なって座ってもよい」という数学的な許可証が、この「\(-1\)」なのです。

\subsection{ボース・アインシュタイン凝縮（BEC）：巨大なスーパーアトムの誕生}

この性質が極限まで現れるのが、物質を絶対零度近くまで極低温に冷やしたときです。 熱の揺らぎ（バラバラに動こうとする邪魔なエネルギー）がなくなると、無数のボース粒子たちは「みんなで一番エネルギーの低い、全く同じ状態（一番下の席）」に一斉にドッと落ち込みます。

すると驚くべきことが起きます。数万、数億個の粒子が個別の粒としてのアイデンティティ（個性）を完全に失い、互いの波がピタリと重なり合って、全体で\textbf{「一つの巨大なスーパーアトム（超原子）」のように振る舞い始めるのです。これを「ボース・アインシュタイン凝縮（BEC）」}と呼びます。

本来はミクロなスケールでしか見えないはずの「波としての性質（量子力学）」が、何億個もの粒子がほぼ同じ状態にそろうことで、マクロな目に見えるサイズの「一つの巨大な波」にまで拡大されて現れる。これは物理学における究極の現象の一つです。 そして、このボース粒子の「同期して群れる性質」こそが、次章（第12章）で解説する、足並みの揃った光のビーム「レーザー」や、電気抵抗がゼロになる「超伝導」といった、現代社会を支えるテクノロジーの重要な源泉となっていくのです。

\section{統計力学とAIの美しいシンクロ：ボルツマンマシンからソフトマックスへ}

ここまでの物理学の話が、現代のAI（人工知能）とどう関係するのでしょうか？ 実は、現在世界中で動いているAI（画像認識からChatGPTのような大規模言語モデルまで）のネットワークの根幹には、11.1節で紹介した統計力学の数式と概念が、そっくりそのまま直輸入されているのです。

\subsection{AIの源流：「イジングモデル」とボルツマンマシン}

AIの歴史において、物理学とAIが最も直接的かつ劇的に融合したのが、ジェフリー・ヒントンらによって考案された\textbf{「ボルツマンマシン」}です。

その土台となったのは、統計力学における\textbf{「イジングモデル（Ising Model）」}という理論です。これは元々、磁石の性質を説明するための数学モデルでした。物質の中にある無数の原子を「上向き（+1）」か「下向き（-1）」のどちらかの状態をとる小さな磁石（スピン）とみなし、隣り合うスピン同士が「同じ向きになりたがる」あるいは「逆向きになりたがる」という相互作用のエネルギーを計算するものです。

ヒントンらは、AIのニューラルネットワークのニューロン（ノード）が「発火するか（1）しないか（0）」の状態を、この物理的な粒子の「スピン」に見立てました。そして、ニューロン同士の繋がりの強さを「スピン間の相互作用」として、ネットワーク全体に対して物理学と全く同じように\textbf{「エネルギー関数（\(E\)）」}を定義したのです。AIが正解に近い望ましい状態になるほどシステム全体のエネルギーが低く、デタラメな状態だとエネルギーが高くなるように設定されました。

\subsection{シミュレーテッド・アニーリング（焼きなまし法）}

このネットワークを「最もエネルギーが低い状態（最適な答え）」に落ち着かせるために、彼らは物理学の\textbf{「アニーリング（焼きなまし）」}という現象をアルゴリズムとして利用しました。

刀鍛冶が鋼を打って強い日本刀を作るとき、金属を高熱にしてから「ゆっくりと」冷やしていきます。すると、金属内部の原子が熱で激しく揺り動かされながら、徐々に最も無駄のない綺麗な結晶構造（エネルギーの絶対的な底）へとピタリと収まっていきます。 AIの学習においても、仮想的な「温度（\(T\)）」のパラメータを持たせます。最初は高温に設定してランダムに激しく探索させ（局所的な最適解の浅い水たまりから無理やり脱出させ）、そこから徐々に温度を下げていくことで、最終的に究極の最適解に落ち着かせるのです。これを\textbf{「シミュレーテッド・アニーリング（Simulated Annealing）」}と呼びます。

このネットワークがある特定のパターンの状態をとる確率は、11.1節の「ボルツマン分布（\(P \propto e^{-E/kT}\)）」に完全に一致するように作られているため、「ボルツマンマシン」と名付けられました。仮想的な物理現象をコンピュータ内で再現することで、AIは学習という知的作業をやってのけたのです。

\subsection{ソフトマックス関数と分配関数}

この統計力学の遺伝子は、現代の最新AIにも脈々と受け継がれています。 AIが「画像に写っているのは犬か、猫か」を確率で出力したり、ChatGPTが「次にくる単語の確率」を計算するとき、最後に出力をパーセンテージ（0〜100\%）に変換するための関数を\textbf{「ソフトマックス関数（Softmax function）」}と呼びます。

ソフトマックス関数の数式を見てみましょう。ある選択肢 \(i\)（例えば「猫」）のスコアを \(x_i\) としたとき、それが選ばれる確率 \(P(i)\) は次のように計算されます。

\[
P(i) = \frac{e^{x_i}}{\sum_j e^{x_j}}
\]

11.1節のボルツマン分布と分配関数の式を思い出してください。

AIの「スコア \(x_i\)」を物理学の「マイナスのエネルギー（\(-E_i/k_B T\)）」に置き換えると、AIのソフトマックス関数は、統計力学のボルツマン分布と同じ形になります。 そして、ソフトマックス関数の分母（すべてのスコアの指数関数の和）は、統計力学のマスターキーである\textbf{「分配関数（\(Z\)）」そのもの}に対応します。

\subsection{AIにおける「分配関数の呪い」とMCMC}

AIの学習において、最も計算が難しく、コンピュータを絶望させるのは、この「分配関数（分母）」の計算です。 選択肢が「犬か猫か（2通り）」なら簡単ですが、AIが複雑な画像や文章を生成する「エネルギーベースモデル」などの場合、あり得るすべての画像のパターン（何億×何億通り）の分母を計算しなければならず、ここでも「次元の呪い」が発生します。

この計算爆発を回避するために、現代のAI研究者たちは、皮肉なことに再び「統計物理学」の手法に頼ることになりました。それが\textbf{「マルコフ連鎖モンテカルロ法（MCMC）」}です。 AIはすべてのパターンをバカ正直に足し算するのを諦め、ランダムに生成したサンプル（熱的な揺らぎ）を使って、統計力学的に「分配関数の近似値」を賢く推測しながら学習を進めているのです。

物理学者が \ensuremath{10^{23}} 個の粒子のカオスを計算するために編み出した「統計力学（ボルツマン分布と分配関数）」は、100年の時を超えて、AIが何十億個のパラメータの中から「確率の高い答え」を導き出すための数学基盤として、現代のデジタル空間で脈々と生き続けています。

そして次章では、この量子統計（フェルミ粒子とボース粒子）のルールが、いかにして私たちの社会を根底から支える「半導体」や「超伝導」といったテクノロジーとして創発していくのかを見ていきます。

\section*{【第11章 章末ディスカッション課題】}

「粒子に個性がなく、完全に区別できない」というミクロの性質が、マクロな世界で「ペチャンコに潰れない物質（フェルミの海）」や「巨大な一つの波（BEC）」といった劇的な現象（創発）を生み出しました。もし人間社会において「個人の区別が全くなくなった」としたら、どのような「マクロな社会現象」が起きると思いますか？ 物理学のアナロジーを使って自由に想像してみましょう。

AIが「次にくる単語」や「画像の意味」を選ぶためのソフトマックス関数が、大自然の「熱エネルギーの落ち着く先」を決めるボルツマン分布と全く同じ数式であることは、単なる偶然でしょうか？ 物理学の法則と情報科学（AI）のアルゴリズムが、深いレベルで一致してしまう理由について、あなたの考えを議論してみてください。


\section*{参考文献（学術誌）}
\begin{enumerate}[leftmargin=2.4em]
\item Bose, S. N. ``Plancks Gesetz und Lichtquantenhypothese.'' \textit{Zeitschrift fur Physik}, 26, 178--181, 1924. DOI: \url{https://doi.org/10.1007/BF01327326}.
\item Fermi, E. ``Zur Quantelung des idealen einatomigen Gases.'' \textit{Zeitschrift fur Physik}, 36, 902--912, 1926. DOI: \url{https://doi.org/10.1007/BF01400221}.
\item Dirac, P. A. M. ``On the theory of quantum mechanics.'' \textit{Proceedings of the Royal Society A}, 112(762), 661--677, 1926. DOI: \url{https://doi.org/10.1098/rspa.1926.0133}.
\item Jaynes, E. T. ``Information theory and statistical mechanics.'' \textit{Physical Review}, 106(4), 620--630, 1957. DOI: \url{https://doi.org/10.1103/PhysRev.106.620}.
\end{enumerate}


